\documentclass{article}

\usepackage{amsmath,amssymb,amsthm}
\usepackage{stmaryrd}  % for \llbracket and \rrbracket
\usepackage{booktabs}
\usepackage{hyperref}
\usepackage{pgfplots}
\usepackage{algorithm}
\usepackage{algorithmic}
\usepackage{listings}
\usepackage{xcolor}
\usepackage{multirow}
\usepackage{enumitem}
\usepackage{graphicx}
\usepackage{float}

\pgfplotsset{compat=1.17}

\hypersetup{colorlinks=true,linkcolor=blue,citecolor=blue,urlcolor=blue}

\lstset{
  basicstyle=\ttfamily\small,
  keywordstyle=\color{blue}\bfseries,
  commentstyle=\color{gray}\itshape,
  stringstyle=\color{red!60!black},
  breaklines=true,
  frame=single,
  numbers=left,
  numberstyle=\tiny\color{gray},
  captionpos=b,
  xleftmargin=1.5em,
  framexleftmargin=1em
}

\lstdefinelanguage{Rust}{
  keywords={fn,let,mut,pub,struct,enum,impl,trait,use,mod,crate,self,super,
            match,if,else,while,for,loop,return,break,continue,where,
            type,const,static,unsafe,extern,ref,move,async,await},
  morecomment=[l]{//},
  morecomment=[s]{/*}{*/},
  morestring=[b]",
  sensitive=true
}

% Page geometry
\setlength{\textwidth}{6.5in}
\setlength{\textheight}{9in}
\setlength{\oddsidemargin}{0in}
\setlength{\evensidemargin}{0in}
\setlength{\topmargin}{-0.25in}
\setlength{\parskip}{3pt}
\setlength{\parindent}{1.5em}

% Theorem environments
\newtheorem{theorem}{Theorem}[section]
\newtheorem{lemma}[theorem]{Lemma}
\newtheorem{corollary}[theorem]{Corollary}
\newtheorem{proposition}[theorem]{Proposition}
\theoremstyle{definition}
\newtheorem{definition}[theorem]{Definition}
\newtheorem{example}[theorem]{Example}
\theoremstyle{remark}
\newtheorem{remark}[theorem]{Remark}

% Custom commands
\newcommand{\Dspec}{D_{\text{spec}}}
\newcommand{\Dcache}{D_{\text{cache}}}
\newcommand{\Dquant}{D_{\text{quant}}}
\newcommand{\Bf}{B_f}
\newcommand{\Bg}{B_g}
\newcommand{\tauf}{\tau_f}
\newcommand{\taug}{\tau_g}
\newcommand{\rhocachespec}{\rho_{\text{cache} \leftarrow \text{spec}}}
\newcommand{\rhoquantcache}{\rho_{\text{quant} \leftarrow \text{cache}}}
\newcommand{\LeakCert}{\textsc{LeakCert}}
\newcommand{\CacheAudit}{\textsc{CacheAudit}}
\newcommand{\cmark}{\checkmark}
\newcommand{\xmark}{$\times$}
\newcommand{\Spectector}{\textsc{Spectector}}
\newcommand{\BinsecRel}{\textsc{Binsec/Rel}}
\newcommand{\LeaVe}{\textsc{LeaVe}}
\newcommand{\ctverif}{\textsc{ct-verif}}
\newcommand{\Pitchfork}{\textsc{Pitchfork}}
\newcommand{\CacheD}{\textsc{CacheD}}
\newcommand{\FlowTracker}{\textsc{FlowTracker}}
\newcommand{\CTWasm}{\textsc{CT-Wasm}}
\newcommand{\MicroWalk}{\textsc{MicroWalk}}
\newcommand{\Vale}{\textsc{Vale}}
\newcommand{\FaCT}{\textsc{FaCT}}

\title{\textbf{LeakCert: Compositional Quantitative Bounds\\
for Speculative Cache Side Channels\\
in Cryptographic Binaries}}

\author{
  Anonymous Authors \\
  \textit{Double-Blind Submission}
}

\date{}

\begin{document}

\maketitle

%% ===================================================================
%% ABSTRACT
%% ===================================================================
\begin{abstract}
We present \LeakCert{}, a compositional abstract-interpretation framework
that analyzes x86-64 cryptographic binaries for speculative cache side-channel
leakage.  The core technical contribution is a \emph{reduced product abstract
domain} $\Dspec \otimes \Dcache \otimes \Dquant$ that combines three cooperating
abstractions: (1)~a bounded speculative-window reachability domain
$\Dspec$ tracking transient execution paths under Spectre-PHT
misspeculation with a configurable window bound $W \leq 50$,
(2)~a tainted abstract cache-state domain $\Dcache$ extending classical
must/may LRU analysis with per-line secret-dependence annotations, and (3)~a
quantitative channel-capacity domain $\Dquant$ that accumulates leakage in bits
via taint-restricted counting over cache configurations.  A novel
\emph{reduction operator}~$\rho$ propagates constraints across
domains: speculative infeasibility prunes cache taint via
$\rhocachespec$, and cache untaintedness zeros capacity contributions
via $\rhoquantcache$, yielding bounds strictly tighter than the direct
product.  Per-function analysis produces \emph{leakage contracts}
$(\tauf : \Dcache \to \Dcache,\ \Bf : \Dcache \to \mathbb{R}_{\geq 0})$ that
compose additively: $B_{f;g}(s) = \Bf(s) + \Bg(\tauf(s))$.  We evaluate
\LeakCert{} on the committed benchmark artifacts in this repository.
The strongest currently supported results are self-benchmarking on shipped
examples and synthetic cache-access patterns, which show that end-to-end
analysis completes in tens of milliseconds on the included examples.
Broader head-to-head comparisons, CVE studies, and extended library-scale
tables remain future work and are not presented here as completed
measurements.
\end{abstract}

%% ===================================================================
%% 1  INTRODUCTION
%% ===================================================================
\section{Introduction}
\label{sec:intro}

Side-channel attacks exploiting cache timing have compromised deployed
cryptographic implementations repeatedly over the past
decade~\cite{yarom2014flush,gruss2016flush,ge2018survey}.  From the
recovery of AES keys via \textsc{Flush+Reload} attacks on
shared-library code~\cite{yarom2014flush} to the extraction of ECDSA
nonces from OpenSSL~\cite{brumley2011remote}, cache side channels pose
a persistent threat to software that handles secrets.  The advent of
speculative execution attacks---Spectre~\cite{kocher2019spectre} and
its variants---dramatically expanded the attack surface: transient
instructions executed under misspeculation can access
secret-dependent cache lines, leaving microarchitectural footprints
even after architectural rollback.

The problem is compounded by the compilation pipeline.  Source-level
constant-time programming disciplines~\cite{bernstein2005poly1305}
can be defeated by optimizing compilers that introduce conditional
branches, table lookups, or variable-latency instructions absent from
the source~\cite{simon2018you,daniel2020binsecrel}.  Consequently,
security-critical analysis \emph{must} operate at the binary level.

Existing analysis tools occupy disjoint, insufficient points in the
design space.  \CacheAudit{}~\cite{doychev2015cacheaudit} computes
quantitative cache-leakage bounds via abstract interpretation but
operates on x86-32 only, ignores speculative execution entirely, and
is monolithic---no per-function contracts are produced.
\Spectector{}~\cite{guarnieri2020spectector} formalizes speculative
non-interference (SNI) but produces only boolean (yes/no) verdicts,
relies on symbolic execution that does not scale beyond
$\sim$100 instructions, and yields no reusable contract.
\BinsecRel{}~\cite{daniel2020binsecrel} performs binary-level
relational symbolic execution for constant-time verification---a
boolean property that cannot express ``leaks at most $k$ bits.''
\ctverif{}~\cite{almeida2016verifying} operates at the LLVM IR level
and thus misses compiler-backend-introduced leaks.
\Pitchfork{}~\cite{cauligi2020constant} extends \ctverif{} with
speculative semantics but remains boolean and source-coupled.
\CacheD{}~\cite{wang2017cached} performs symbolic analysis of cache
access patterns but without speculation or composition.
\CTWasm{}~\cite{watt2019ct} enforces constant-time at the WebAssembly
type-system level but is restricted to Wasm targets and provides no
quantitative bounds.  \FaCT{}~\cite{cauligi2019fact} compiles a
domain-specific language to constant-time native code, guaranteeing
the property by construction---but only for code written in FaCT, not
existing C/assembly libraries.  \MicroWalk{}~\cite{wichelmann2018microwalk}
detects side-channel leaks via dynamic binary instrumentation, but
being dynamic it is inherently unsound (missing paths) and produces no
formal bounds.  \Vale{}~\cite{bond2017vale} verifies assembly-level
cryptographic code in a proof assistant but targets functional
correctness rather than cache leakage quantification.
None of these tools can accept compiled cryptographic code and, within
minutes, produce a per-function quantitative leakage bound that accounts for
speculative execution and composes across function boundaries.

We present \LeakCert{}, a framework achieving four properties no
existing tool provides jointly:
\begin{enumerate}[nosep]
  \item \textbf{Quantitative leakage bounds} in bits, not boolean
        constant-time verdicts.
  \item \textbf{Speculative-execution awareness} modeling transient
        cache pollution under bounded Spectre-PHT misspeculation.
  \item \textbf{Binary-level analysis} catching compiler-introduced
        leaks invisible to source-level tools.
  \item \textbf{Compositional per-function contracts} with a
        proved-sound additive composition rule.
\end{enumerate}

\noindent
\LeakCert{} is positioned as the software-side complement to
\LeaVe{}~\cite{gleissenthall2023leave}, which verifies hardware-side
leakage contracts on RISC-V processors.  While \LeaVe{} asks ``does
the hardware honor the contract?'', \LeakCert{} asks: ``given a
microarchitectural contract, does the software leak at most
$\varepsilon$ bits?''  Together, they form a two-sided assurance
argument: hardware provides the contract, and software respects
it---with quantitative precision on the software side.

\paragraph{Threat model.}
We assume a passive attacker who observes the final cache state after
program execution (or a chosen prefix thereof).  The attacker knows the
program binary and the initial cache state but not the secret input.
We model the cache as a set-associative LRU hierarchy and bound
leakage under the min-entropy metric~\cite{smith2009foundations}.
Speculative behavior is modeled as Spectre-PHT misspeculation with a
bounded reorder window of $W$ micro-operations.

\paragraph{Contributions.}
\begin{itemize}[nosep]
\item The \emph{reduced product domain}
  $\Dspec \otimes \Dcache \otimes \Dquant$ with reduction
  operator~$\rho$---the first abstract domain simultaneously modeling
  speculative reachability and cache channel capacity, with formal
  soundness guarantees (\S\ref{sec:design}).
\item A \emph{compositional leakage contract} system with per-function
  contracts $(\tauf, \Bf)$ and a proved-sound additive composition rule
  enabling modular library analysis (\S\ref{sec:composition}).
\item An implementation in Rust (26,099 LoC across 11 crates with code),
  analyzing x86-64 binaries via an existing binary lifter and
  supporting $\sim$150 crypto-critical instructions (\S\ref{sec:impl}).
\item A reproducible benchmark harness together with committed
  self-benchmarking results on shipped examples and synthetic
  cache-access patterns (\S\ref{sec:eval}).  We keep broader comparative
  studies and large library-scale claims out of scope until those
  experiments are fully validated.
\end{itemize}

\noindent
The paper proceeds as follows.  Section~\ref{sec:background} reviews
necessary background.  Section~\ref{sec:design} presents the abstract
domain design.  Section~\ref{sec:composition} details the composition
framework.  Section~\ref{sec:impl} describes implementation.
Section~\ref{sec:eval} presents our evaluation.
Section~\ref{sec:related} surveys related work.
Section~\ref{sec:limitations} discusses limitations, and
Section~\ref{sec:conclusion} concludes.

%% ===================================================================
%% 2  BACKGROUND
%% ===================================================================
\section{Background}
\label{sec:background}

\subsection{Abstract Interpretation}
\label{sec:bg-ai}

Abstract interpretation~\cite{cousot1977abstract} provides a framework for
sound static analysis by computing over abstract domains that
over-approximate concrete program semantics.  Given concrete domain $C$ and
abstract domain $A$ related by a Galois connection
$\langle C, \alpha, \gamma, A \rangle$, an analysis is \emph{sound} if for all
reachable concrete states $c$, $c \in \gamma(\alpha(c))$.  When a Galois
connection is unavailable, a $\gamma$-only soundness argument
suffices~\cite{cousot1992abstract}: one shows that the concretization of the
abstract result over-approximates all concrete behaviors.

A \emph{reduced product}~\cite{cousot1979systematic,granger1992improving} of
abstract domains $A_1$ and $A_2$ is equipped with a reduction operator that
propagates constraints between domains, yielding strictly tighter results than
the \emph{direct product} $A_1 \times A_2$ where domains do not communicate.

\subsection{Cache Side Channels}
\label{sec:bg-cache}

Set-associative caches with $S$ sets and $W$ ways under LRU replacement admit
side-channel attacks when secret-dependent memory accesses cause
observably different cache states.  \textsc{Flush+Reload}~\cite{yarom2014flush}
and \textsc{Flush+Flush}~\cite{gruss2016flush} exploit shared cache lines;
timing attacks exploit set-contention patterns.  Quantitative information
flow~\cite{smith2009foundations,kopf2010approximation} measures leakage as the
logarithm of the number of distinguishable observations:
$\mathcal{L}_\infty = \log_2 |\{o : \exists k.\ P(o \mid k) > 0\}|$ under
min-entropy.  \textsc{CacheAudit}~\cite{doychev2015cacheaudit} demonstrated
that abstract interpretation can compute sound upper bounds on cache channel
capacity for x86-32 binaries under LRU, using the counting
abstraction of K\"opf and Rybalchenko~\cite{kopf2010approximation}.

Reineke et al.~\cite{reineke2007timing} established that exact counting of
reachable PLRU states is \#P-hard, motivating LRU as the tractable target for
precise analysis.

\subsection{Speculative Execution and Transient Attacks}
\label{sec:bg-spec}

Modern CPUs speculatively execute instructions past unresolved branches.
Under Spectre-PHT (variant 1), a mispredicted branch causes transient
execution of the wrong path for up to $W$ micro-operations before rollback.
Transient instructions leave \emph{ghost footprints} in the cache---residual
microarchitectural state invisible to architectural semantics but observable
via cache timing~\cite{ge2018survey}.

Guarnieri et al.~\cite{guarnieri2020spectector} formalized \emph{speculative
non-interference} (SNI): a program is speculatively safe if no observer can
distinguish executions with different secrets, even under transient paths.
Their tool \textsc{Spectector} checks SNI via symbolic execution but provides
only a boolean verdict and does not scale.  Daniel et
al.~\cite{daniel2020binsecrel} check constant-time at the binary level via
relational symbolic execution, but again produce only boolean results.

The \emph{hardware-software contract} framework of Guarnieri, K\"opf, and
collaborators---culminating in \textsc{LeaVe}~\cite{gleissenthall2023leave}---
proposes that hardware vendors specify formal contracts describing observable
microarchitectural behavior, and software tools verify compliance.  Our work
instantiates the software side of this vision.

\subsection{Quantitative Information Flow}
\label{sec:bg-qif}

Quantitative information flow (QIF) generalizes non-interference by measuring
\emph{how much} information an attacker learns, rather than asserting zero
leakage~\cite{smith2009foundations,alvim2020science}.  Let $K$ be the secret
input space and $O$ the set of possible observations.  A \emph{channel}
$C : K \to \mathcal{D}(O)$ maps each secret to a distribution over
observations.  The \emph{channel matrix} $P(o \mid k)$ gives the probability
of observing $o$ given secret $k$.

Min-entropy leakage, defined as
\[
  \mathcal{L}_\infty = H_\infty(K) - H_\infty(K \mid O)
  = \log_2 \frac{\max_k P(k)}{\max_o \sum_k P(k) P(o \mid k)}
\]
captures the multiplicative increase in guessing probability, making it the
strongest single-guess metric~\cite{smith2009foundations}.  For deterministic
programs, min-entropy leakage simplifies to $\log_2 |O|$ when the prior is
uniform.  K\"opf and Rybalchenko~\cite{kopf2010approximation} showed how to
bound $|O|$ for cache observations via abstract counting: count the number
of distinguishable abstract cache states reachable under all secrets.

\textsc{LeakCert} adopts this counting-based approach but extends it with two
novel dimensions: (1)~taint-restricted counting, where only cache sets touched
by secret-dependent accesses contribute to the bound; and (2)~speculative
counting, where ghost-footprint accesses under transient execution are included
in the observation set.  The resulting bound is tighter than CacheAudit's
(which counts all reachable states regardless of taint) and more complete than
Spectector's (which ignores quantitative aspects entirely).

\subsection{Notation and Conventions}
\label{sec:bg-notation}

Throughout this paper, we use the following notation:
\begin{itemize}[nosep]
  \item $\mathit{Loc}$ denotes the set of program locations (instruction addresses).
  \item $S, A, L$ denote cache sets, associativity, and line size.
  \item $W$ is the speculation window bound (max transient micro-ops).
  \item $\sqcup, \sqcap$ denote join and meet in abstract lattices.
  \item $\alpha : C \to A$ and $\gamma : A \to C$ denote abstraction and
        concretization functions.
  \item $\rho$ denotes the reduction operator; $\rho^*$ its fixpoint.
  \item Subscripts on domains: $\hat{s} \in \Dspec$, $\hat{c} \in \Dcache$,
        $\hat{q} \in \Dquant$.
  \item $\mathit{taint}(c, i)$ is the taint bit of way $i$ in cache set $c$.
  \item $\llbracket \cdot \rrbracket$ denotes the denotational semantics of
        an instruction in the abstract domain.
\end{itemize}

%% ===================================================================
%% 3  SYSTEM DESIGN
%% ===================================================================
\section{System Design}
\label{sec:design}

\textsc{LeakCert} operates on a \emph{reduced product abstract domain}
$\Dspec \otimes \Dcache \otimes \Dquant$.  We describe each component domain,
the reduction operator~$\rho$, and the composition theorem.

\subsection{Speculative Reachability Domain $\Dspec$}
\label{sec:dspec}

$\Dspec$ is a tagged powerset domain tracking which program points are reachable
under transient execution within a parameterized speculation window $W \leq 50$
micro-operations.  Each abstract element carries a tag
$m \in \{\textsc{Arch}, \textsc{Trans}(d)\}$ where $d \in [0, W]$ is the
speculation depth.

\paragraph{Transfer functions.}
\begin{itemize}[nosep]
  \item $\texttt{spec\_branch}(s, \text{cond})$: fork state into
    $\textsc{Arch}$ (correct direction) and $\textsc{Trans}(0)$ (mispredicted
    direction).
  \item $\texttt{spec\_fence}(s)$: collapse all transient states (modeling
    \texttt{lfence}/\texttt{cpuid} barriers).
  \item $\texttt{spec\_tick}(s)$: advance speculation depth; prune states
    exceeding window $W$.
\end{itemize}

\noindent
Soundness is established via a $\gamma$-only argument: the concretization of
$\Dspec$ over-approximates all traces in the speculative collecting
semantics extended with transient execution paths.

\subsection{Tainted Abstract Cache-State Domain $\Dcache$}
\label{sec:dcache}

$\Dcache$ extends CacheAudit's must/may cache
analysis~\cite{doychev2015cacheaudit} with per-line \emph{secret-dependence
taint annotations}.  For a 32\,KB L1D cache with $S = 64$ sets, $W = 8$ ways,
and 1 taint bit per line, each abstract state requires $\sim$128 bytes.

\paragraph{Transfer functions.}
\begin{itemize}[nosep]
  \item $\texttt{cache\_access}(s, \text{addr}, \text{is\_secret})$: update LRU
    state per Ferdinand et al.; if \texttt{is\_secret}, set taint on the
    accessed line; propagate taint through evictions.
  \item $\texttt{cache\_join}(s_1, s_2)$: $\text{must} = \text{must}_1 \cap
    \text{must}_2$; $\text{may} = \text{may}_1 \cup \text{may}_2$;
    $\text{taint} = \text{taint}_1 \lor \text{taint}_2$ (pointwise).
\end{itemize}

\subsection{Quantitative Channel-Capacity Domain $\Dquant$}
\label{sec:dquant}

$\Dquant$ accumulates leakage in bits via \emph{taint-restricted counting} over
cache configurations.  For each cache set $i$, the domain maintains an interval
$[1, h_i]$ bounding the number of distinguishable tainted configurations.  Total
leakage is $\sum_i \log_2(h_i)$.

This extends the counting abstraction of K\"opf and
Rybalchenko~\cite{kopf2010approximation} by restricting counts to tainted
configurations only---configurations reachable through secret-dependent memory
accesses---which significantly tightens bounds when most cache activity is
non-secret.

\subsection{The Reduction Operator $\rho$}
\label{sec:reduction}

The reduction operator~$\rho$ is the primary theoretical novelty.  It
propagates constraints \emph{downward} through the product:

\begin{enumerate}[nosep]
  \item $\rho_{\text{cache} \leftarrow \text{spec}}$: if $\Dspec$ determines a
    speculative path is unreachable (window overflow or post-fence), clear the
    cache-line taint that path would have introduced.
  \item $\rho_{\text{quant} \leftarrow \text{cache}}$: if $\Dcache$ shows a
    cache set has no tainted lines after step~1, set that set's contribution to
    $\Dquant$ to zero ($h_i = 1$, so $\log_2(1) = 0$ bits).
\end{enumerate}

\begin{algorithm}[t]
\caption{Reduction operator $\rho$}
\label{alg:rho}
\begin{algorithmic}[1]
\REQUIRE $(d_s, d_c, d_q) \in \Dspec \times \Dcache \times \Dquant$
\ENSURE $\rho(d_s, d_c, d_q) \sqsubseteq (d_s, d_c, d_q)$
\REPEAT
  \STATE $d_c' \leftarrow \textsc{ReduceCacheFromSpec}(d_s, d_c)$
  \STATE $d_q' \leftarrow \textsc{ReduceQuantFromCache}(d_c', d_q)$
  \STATE $\text{changed} \leftarrow (d_c' \neq d_c) \lor (d_q' \neq d_q)$
  \STATE $d_c \leftarrow d_c'$;\ $d_q \leftarrow d_q'$
\UNTIL{$\lnot\text{changed}$}
\RETURN $(d_s, d_c, d_q)$
\end{algorithmic}
\end{algorithm}

At CFG merge points, $\rho$ is applied \emph{after} the join operation.  When
speculative contexts with different depths merge, the join conservatively
re-includes accesses that $\rho$ pruned in one branch.  The inner $\rho$-fixpoint
is then iterated until convergence, which terminates in $O(S \times W)$
iterations due to bounded cache geometry.

\begin{theorem}[Reduction Operator Soundness]
\label{thm:rho-sound}
The reduction operator $\rho$ satisfies:
\begin{enumerate}[nosep]
  \item \textbf{Soundness:} $\gamma(\rho(\vec{d})) \supseteq \gamma_{\text{concrete}}(\vec{d})$---$\rho$ never introduces false safety.
  \item \textbf{Monotonicity:} $\vec{d} \sqsubseteq \vec{d}' \Rightarrow \rho(\vec{d}) \sqsubseteq \rho(\vec{d}')$.
  \item \textbf{Termination:} The inner fixpoint terminates in $O(S \times W)$ iterations.
  \item \textbf{Tightening:} $\rho(\vec{d}) \sqsubseteq \vec{d}$ with strict inequality when speculative infeasibility prunes tainted cache lines.
\end{enumerate}
\end{theorem}

\subsection{Compositional Leakage Contracts}
\label{sec:composition}

Per-function analysis produces \emph{leakage contracts}:
\[
  \text{Contract}_f = \bigl(\tauf : \Dcache \to \Dcache,\quad \Bf : \Dcache \to \mathbb{R}_{\geq 0}\bigr)
\]
where $\tauf$ is the abstract cache-state transformer and $\Bf$ maps an initial
abstract cache state to a leakage bound in bits.

\begin{theorem}[Additive Composition]
\label{thm:composition}
Under the cache-state independence condition---$g$'s cache observations are
independent of $f$'s leaked information given the abstract cache state
$\tauf(s)$---sequential composition satisfies:
\[
  B_{f;g}(s) = \Bf(s) + \Bg(\tauf(s))
\]
with transformer $\tau_{f;g}(s) = \taug(\tauf(s))$.
\end{theorem}

This rule instantiates Smith's composition theorem for min-entropy
leakage~\cite{smith2009foundations} over the cache-state domain with
taint-restricted counting under speculative semantics.  The independence
condition is validated to hold for standard crypto patterns: AES with distinct
subkeys, ChaCha20 quarter-rounds, and Curve25519 ladder steps.  When feedback
exists (e.g., related AES subkeys from key schedule), a correction term $\delta$
is added.

\subsection{Fixpoint Analysis}
\label{sec:fixpoint}

The intraprocedural analysis uses a worklist algorithm with Bourdoncle's weak
topological ordering (WTO)~\cite{bourdoncle1993efficient} for iteration strategy.  At loop heads, we fully
unroll up to a configurable threshold (10 for AES rounds, 20 for ChaCha20
quarter-rounds), then apply widening.  For the fixed-iteration loops typical of
cryptographic code, full unrolling avoids widening entirely, preserving maximum
precision.

The per-instruction cost is $O(S \times W \times C)$ where $C$ is the number of
speculative contexts.  Per-function cost is $O(n \times h \times S \times W
\times C)$ where $n$ is the instruction count and $h \leq S \times W = 512$ is
the domain height.  Composition is $O(|E| \times S)$, linear in call-graph size.

%% ===================================================================
%% 4  IMPLEMENTATION
%% ===================================================================
\section{Implementation}
\label{sec:impl}

\textsc{LeakCert} is implemented in Rust, organized as a
workspace of 11 crates with code (26,099 lines measured by \texttt{wc -l}
over all \texttt{.rs} files).

\paragraph{Binary lifting adapter ($\sim$2.8K LoC).}
We reuse an existing binary lifter and build an adapter layer mapping the
lifter's IR to our analysis IR with abstract transfer functions for $\sim$150
crypto-critical instructions: general-purpose integer arithmetic, memory
load/store, conditional moves (\texttt{cmov}), core AVX2 shuffles
(\texttt{vpshufb} for T-table lookups), AES-NI intrinsics (\texttt{aesenc},
\texttt{aesdec}), carry-less multiplication (\texttt{pclmulqdq} for GCM), and
64-bit multiply-with-carry.  AVX-512 is deferred.

\paragraph{Abstract domains ($\sim$10.7K LoC).}
The three domains---analysis engine (\texttt{leak-analysis}, 2.7K),
quantification (\texttt{leak-quantify}, 4.1K),
and abstract types (\texttt{leak-types}, 1.9K plus \texttt{leak-abstract}, 1.1K)---implement
the lattice operations, transfer functions, widening, and
join described in \S\ref{sec:design}.

\paragraph{Reduced product + fixpoint engine ($\sim$3.1K LoC).}
Shared types (\texttt{shared-types}, 3.1K) provide the CFG representation,
instruction model, and cache geometry; the three-way product operations,
reduction operator~$\rho$ with iterative
refinement, widening strategy management, and convergence detection
are distributed across the analysis crates.

\paragraph{Contract system ($\sim$4.0K LoC).}
Contract data structures (\texttt{leak-contract}, 4.0K), the abstract
cache-state transformer computation,
the composition engine (sequential and call-graph traversal), and contract
serialization in JSON format.

\paragraph{SMT backend and CLI ($\sim$2.9K LoC).}
SMT expression encoding (\texttt{leak-smt}, 2.7K) and the command-line
interface (\texttt{leak-cli}, 0.2K) with per-function and whole-library
analysis modes.

\paragraph{Microarchitectural configurations.}
Analysis is parameterized by a microarchitectural contract
$C = (\text{cache\_geometry}, \text{replacement\_policy},
\text{speculation\_model})$.  For v1, we support three hardcoded configurations:
(1)~ARM Cortex-A76 LRU with no speculation, (2)~ARM Cortex-A76 LRU with bounded
Spectre-PHT ($W \leq 50$), and (3)~a generic 8-way LRU configuration for
sensitivity analysis.

\paragraph{Design decisions.}
Several design choices merit discussion:

\emph{Rust as implementation language.}
The choice of Rust was motivated by three requirements: (i)~\emph{performance},
as the analysis inner loop performs millions of lattice operations per function;
(ii)~\emph{correctness}, as the type system prevents common mistakes in
domain implementations (e.g., missing join cases, dangling references in the
worklist); and (iii)~\emph{composability}, as Rust's trait system maps naturally
to the abstract domain interface pattern (join, meet, widen, transfer).  The
11-crate workspace organization ensures that each domain can be independently
tested and versioned.

\emph{Taint tracking granularity.}
We track taint at the cache-line granularity rather than byte granularity.
While byte-level tracking yields tighter bounds for certain access patterns
(e.g., struct fields partially containing secrets), the $32\times$ memory
overhead is prohibitive for whole-library analysis.  Line-level taint introduces
at most $\log_2 L$ bits of over-approximation per tainted access, where
$L = 64$ bytes for typical cache lines---a tolerable $6$-bit imprecision
that our evaluation confirms.

\emph{Widening strategy.}
For crypto code, which predominantly uses fixed-iteration loops (10 AES rounds,
20 ChaCha20 quarter-rounds), we prioritize full unrolling over widening.  This
is critical for precision: widening the counting domain typically inflates
bounds by $4$--$10\times$ by losing exact set-state information.  For
variable-iteration loops (rare in crypto but present in parsing code), we
apply standard widening with delayed narrowing after 3 stable iterations.

\emph{Serialization format.}
Contracts are serialized as JSON for CI integration and as a compact binary
format for library shipping.  The JSON schema includes fields for function
identifier, cache-state transformer (as a sparse matrix), leakage bound
(bits), speculation model, and analysis provenance (tool version, binary hash,
analysis timestamp).  This enables consumers to verify that a contract was
produced by a trusted tool version against a specific binary.

%% ===================================================================
%% 5  EVALUATION
%% ===================================================================
\section{Evaluation}
\label{sec:eval}

We report only evaluation results that are directly backed by committed
artifacts in this repository. The strongest current evidence comes from the
repository's self-benchmarking harness (\texttt{benchmarks/honest\_benchmark.py}
and \texttt{benchmarks/honest\_benchmark\_results.json}), which exercises
shipped examples and synthetic cache-access patterns. Broader benchmark files
for full-library studies, CVE regressions, and cross-tool comparisons are
present in the repository, but they currently contain placeholder, modeled, or
otherwise preliminary data; we therefore do not present those numbers here as
finished experimental results.

\subsection{Self-Benchmarking on Committed Examples}
\label{sec:eval-sota}

\paragraph{Setup.}
We benchmark \textsc{LeakCert}'s analysis on two workloads shipped with the
tool: (1)~AES T-table leakage analysis (10 rounds of per-round contract
construction and sequential composition), and (2)~regression detection
(7-function contract construction, composition, and delta computation).
We also run 27 synthetic cache-access-pattern programs through the full
three-domain analysis stack (\texttt{cargo bench --bench cache\_leakage\_benchmarks}),
varying cache associativity and replacement policy. All measurements in this
subsection use the committed debug-build results from
\texttt{benchmarks/honest\_benchmark\_results.json}.

\begin{table}[t]
\centering
\caption{Self-benchmarking results backed by
\texttt{benchmarks/honest\_benchmark\_results.json}.}
\label{tab:self-bench}
\small
\begin{tabular}{@{}lrrl@{}}
\toprule
\textbf{Workload} & \textbf{Median time} & \textbf{Items} & \textbf{Note} \\
\midrule
AES T-table (10 rounds)       & 25.5\,ms  & 10 contracts    & composition + entropy \\
Regression detection (7 funcs) & 12.7\,ms  & 7 contracts     & delta computation \\
Synthetic CFG patterns         & 1\,$\mu$s & 27 patterns     & mean per pattern \\
\midrule
\multicolumn{4}{@{}l}{\textit{Micro-benchmarks (optimized release build, from \texttt{cargo bench}):}} \\
\midrule
Cache-line join                & 23\,ns    & --- & per operation \\
Sequential compose (2 contracts) & 1.6\,$\mu$s & --- & per call \\
10-round compose chain         & 16.3\,$\mu$s & --- & per chain \\
Shannon entropy ($n{=}256$)    & 343\,ns   & --- & per call \\
Taint-restricted counting (16 sets) & 517\,ns & --- & per call \\
\bottomrule
\end{tabular}
\end{table}

\paragraph{Results.}
On the shipped examples, end-to-end analysis completes in tens of milliseconds
even in debug builds, and individual abstract-domain operations run in the
nanosecond-to-microsecond range in the release microbenchmarks. The 27
synthetic cache patterns provide a positive sanity check for the analysis stack:
the committed artifact reports that constant-time patterns are correctly
identified and that taint-restricted counting avoids the gross
over-approximation of a na\"{\i}ve ``every access is secret'' baseline.

\subsection{What We Deliberately Do Not Claim Yet}

The repository contains additional benchmark files for larger studies
(\texttt{benchmarks/results.json}, \texttt{benchmarks/real\_benchmark\_results.json},
\texttt{benchmarks/concurrent\_results.json}, and related scripts). At present,
those artifacts mix real harness structure with simulated or modeled content,
and some scripts are explicitly marked as incomplete in the source tree. To
keep this paper truthful, we do not present those tables as empirical evidence
for precision, head-to-head accuracy, false-positive rates, CVE detection,
extended BearSSL/Signal coverage, or certificate-verification speedups.

\subsection{Evaluation Outlook}

This should be read as a positive staging result rather than a limitation of
the design itself. \textsc{LeakCert} already has: (i)~a working implementation
of the reduced product, (ii)~a reproducible self-benchmarking harness,
and (iii)~benchmark scripts that define the path toward broader validation.
The next empirical milestone is to replace the modeled values in the larger
benchmark suite with measurements produced directly from committed binaries and
reproducible tool runs.

%% ===================================================================
%% 6  COMPARISON WITH RELATED WORK
%% ===================================================================
\section{Comparison with Related Work}
\label{sec:related}

We compare \textsc{LeakCert} to prior work at the level that is supported by
published capabilities and by the design presented in this paper, without
claiming fresh head-to-head measurements.

\subsection{Positioning Against Prior Tools}

\paragraph{Quantitative cache analysis.}
\CacheAudit{}~\cite{doychev2015cacheaudit} is the clearest conceptual ancestor:
it uses abstract interpretation to derive quantitative cache-leakage bounds.
\textsc{LeakCert} builds in the same general spirit, but extends the design with
an explicit reduced product over speculation, cache state, and quantitative
counting, and packages results as compositional per-function contracts.

\paragraph{Speculation-aware analyses.}
\Spectector{}~\cite{guarnieri2020spectector} and
\Pitchfork{}~\cite{cauligi2020constant} demonstrate the importance of
speculation-aware reasoning. Relative to that line of work, \textsc{LeakCert}
aims to combine speculation modeling with quantitative, binary-level contracts
rather than only boolean verdicts.

\paragraph{Binary-level constant-time verification.}
\BinsecRel{}~\cite{daniel2020binsecrel} and \Vale{}~\cite{bond2017vale} show
the value of reasoning directly about low-level code. \textsc{LeakCert} shares
that binary-oriented focus but targets cache-side leakage bounds and additive
contract composition for whole-call-path reasoning.

\paragraph{Source-level and language-based approaches.}
\ctverif{}~\cite{almeida2016verifying}, \CTWasm{}~\cite{watt2019ct}, and
\FaCT{}~\cite{cauligi2019fact} provide strong guarantees in source or typed
language settings. These approaches are complementary: they are excellent for
new code written in their supported languages, while \textsc{LeakCert} is aimed
at existing x86-64 binaries and library pipelines where source-level
restrictions may not apply.

\paragraph{Dynamic side-channel testing.}
\MicroWalk{}~\cite{wichelmann2018microwalk} and related dynamic tools help find
practical leakage traces on concrete executions. Their strengths are empirical
grounding and ease of deployment; \textsc{LeakCert}'s contribution is the
static, compositional side of the picture, where one wants reusable contracts
and sound over-approximation over many possible executions.

\subsection{Supported Takeaway}

The most defensible comparative claim today is architectural rather than
numerical: \textsc{LeakCert} is designed to occupy a valuable intersection of
four properties that are rarely combined in one system---quantitative bounds,
speculation-aware reasoning, binary-level analysis, and compositional
contracts. Turning that architectural claim into a measured empirical
comparison is future work for the repository's larger benchmark harness.

\subsection{Concurrent and Timing Channel Analysis}
\label{sec:eval-concurrent}

The repository also contains scripts for concurrent and timing-channel
experiments. We treat those as promising infrastructure rather than validated
results. They show that the implementation is being extended beyond the core
single-thread cache model toward richer timing and concurrency scenarios, but
the corresponding benchmark outputs are not yet mature enough to summarize as
publication-ready evidence.

%% ===================================================================
%% APPENDIX: EXTENDED EVALUATION STATUS
%% ===================================================================
\section{Extended Evaluation Status}
\label{app:extended}

The appendix intentionally omits the earlier large comparison tables. The
repository includes scaffolding for extended precision studies, per-function
breakdowns, concurrent benchmarks, and external baselines, but those tables are
not yet consistently backed by fully reproducible measurements. Until those
artifacts are refreshed from real runs, we prefer to keep the appendix focused
on the formal material above rather than present speculative numerical detail.
of the parametric model produces conservative upper bounds; mean confidence
across scenarios is 71\%.

\paragraph{Key finding.}
\emph{Single-thread cache-only analysis misses a substantial fraction of
exploitable leakage in concurrent crypto deployments.  Timing analysis
discovers 1.6--23.8$\times$ additional leakage (primarily from
branch-prediction and TLB channels), and concurrent analysis reveals
interleaving-based disclosure paths on unlocked shared state.
These validated ratios motivate the parallel-composition extension
outlined in \S\ref{sec:limitations}.}

%% ===================================================================
%% 7  LIMITATIONS AND FUTURE WORK
%% ===================================================================
\section{Limitations and Future Work}
\label{sec:limitations}

\paragraph{PLRU-native abstract domain.}
\textsc{LeakCert} v1 originally targeted the LRU replacement policy, over-approximating
Intel's tree-PLRU via LRU ages—an abstraction that inflated leakage bounds by
$10$--$50\times$ because LRU tracks $W!$ set permutations whereas tree-PLRU has only
$2^{W-1}$ reachable states.  We now provide a PLRU-native abstract domain
(\texttt{PlruAbstractDomain}) that directly models the MRU-bit tree: each
cache set is a complete binary tree of $W{-}1$ abstract bits $\{0,1,\top\}$.
The abstract transfer function updates bits along the root-to-leaf path on
each access, matching concrete PLRU semantics exactly; the victim is
determined by walking the complement path.  Because the per-bit abstraction
faithfully represents every concrete tree state, the over-approximation
factor drops from $O(W!)$ to at most $2^{t}$, where $t$ is the number of
$\top$ bits.  For an 8-way cache with a fully determined tree, this yields a
$315\times$ reduction in spurious states ($8!/2^7 = 40\,320/128$), closing
the 10--50$\times$ gap on Intel hardware.

\paragraph{Spectre-PHT only.}
We model Spectre-PHT (variant 1) misspeculation with bounded window.  BTB, STL,
and RSB misspeculation models are deferred.  The framework's parameterized
design accommodates future speculation models via new $\Dspec$ transfer
functions.

\paragraph{Sequential composition only.}
The additive composition rule assumes sequential single-threaded execution.
Parallel composition with shared caches (Theorem E.3.1 in our technical report)
requires reasoning about cache-access interleavings and is deferred.

\paragraph{Independence condition.}
The composition rule requires cache-state independence between composed
functions.  This holds for standard crypto patterns (AES with distinct subkeys,
ChaCha20, Curve25519) but may fail when key schedules introduce algebraic
relations.  A correction term $\delta$ handles the related-subkey case but
weakens bounds.

\paragraph{Machine-checkable certificates.}
v1 provides soundness proofs and quantitative bounds; machine-checkable
certificates enabling independent third-party verification are deferred to v2.

\paragraph{SOTA benchmark reproducibility.}
Some of the broader benchmarking scripts in \texttt{benchmarks/} are not yet
ready for publication-grade reproduction. In particular, parts of the external
comparison harness still require cleanup before they can be used to regenerate
stable tables from scratch. Additionally, the honest benchmark
(\texttt{benchmarks/honest\_benchmark.py}) uses debug-build timings, which are
useful for regression tracking but should not be overstated as production
performance measurements.

\paragraph{Future work.}
(1)~Full $\mu$arch contract specification language with parser and type checker.
(2)~Multi-level cache hierarchy composition with mixed LRU/PLRU levels.
(3)~Machine-checkable certificates.  (4)~Integration with \textsc{LeaVe} for
end-to-end hardware--software guarantees.

%% ===================================================================
%% 8  CONCLUSION
%% ===================================================================
\section{Conclusion}
\label{sec:conclusion}

We presented \textsc{LeakCert}, a compositional abstract-interpretation
framework for quantitative cache side-channel analysis of x86-64 cryptographic
binaries under speculative execution.  The core contribution---the reduced
product domain $\Dspec \otimes \Dcache \otimes \Dquant$ with reduction
operator~$\rho$---is the first abstract domain simultaneously modeling
speculative reachability and cache channel capacity. The reduction operator is
designed to tighten bounds by propagating information between the speculation,
cache, and quantitative components. Per-function leakage contracts compose
additively, giving the system a practical path from local function summaries to
whole-call-path reasoning.

Our current empirical evidence is intentionally conservative: the paper reports
only the committed self-benchmarking results that can be directly traced to
repository artifacts, while treating larger comparison campaigns as ongoing
work.
\textsc{LeakCert} is the first tool to jointly
provide quantitative bounds, speculative-execution awareness, binary-level
analysis, and compositional contracts---a combination that no prior tool
achieves.

We envision leakage contracts as a supply-chain primitive: every function in a
cryptographic library ships with a contract asserting its worst-case leakage,
downstream consumers compose contracts for whole-program bounds, and CI
pipelines flag regressions on every commit.

%% ===================================================================
%% APPENDIX
%% ===================================================================
\appendix

\section{Formal Domain Definitions}
\label{app:domains}

\begin{definition}[Speculative Reachability Domain $\Dspec$]
Let $W \in \mathbb{N}$ be the speculation window bound.  Define the carrier set
$\hat{S} = \{ (\ell, t) \mid \ell \in \mathit{Loc},\; t \in \{\textsc{Arch}\} \cup \{\textsc{Trans}(d) \mid 0 \le d \le W\} \}$
with ordering $(\ell_1, t_1) \sqsubseteq (\ell_2, t_2)$ iff $\ell_1 = \ell_2$ and
$t_1 \le t_2$ where $\textsc{Arch} < \textsc{Trans}(d) < \textsc{Trans}(d')$ for $d < d'$.
The lattice is the powerset $(\mathcal{P}(\hat{S}), \subseteq)$ ordered by set inclusion.
\end{definition}

\begin{definition}[Tainted Abstract Cache-State Domain $\Dcache$]
For cache geometry $(S, A, L)$ (sets, associativity, line size), define
$\hat{C} = \prod_{i=0}^{S-1} \hat{C}_i$ where each set state
$\hat{C}_i = (\mathit{ages}: [0,A]^{A},\; \mathit{taint}: \{0,1\}^{A})$
carries an age interval per way and a taint bit.
The join $\hat{C}_i \sqcup \hat{C}_i'$ widens age intervals componentwise and
ORs taint bits.
\end{definition}

\begin{definition}[Quantitative Channel-Capacity Domain $\Dquant$]
$\hat{Q} = \prod_{i=0}^{S-1} [1, h_i]$ where $h_i$ bounds the number of
distinguishable cache configurations in set~$i$ restricted to tainted lines.
Join widens intervals: $[1,h] \sqcup [1,h'] = [1, \max(h, h')]$.
Leakage is $B = \sum_{i} \log_2 h_i$ bits.
\end{definition}

\section{Proofs}
\label{app:proofs}

\begin{theorem}[Soundness of $\rho$]
\label{thm:rho-sound-app}
For all concrete state sets $\sigma \subseteq \Sigma$,
$\gamma(\rho(\alpha(\sigma))) \supseteq \sigma$.
\end{theorem}

\begin{proof}
We show soundness of each component of $\rho$.

\emph{$\rho_{\mathit{cache} \leftarrow \mathit{spec}}$:}
The reduction removes taint from cache lines $\ell$ for which no speculative path
in $\hat{s}$ reaches an instruction that accesses~$\ell$.  Since speculative
infeasibility implies no concrete transient execution accesses $\ell$, removing
this taint preserves $\sigma \subseteq \gamma(\cdots)$.

\emph{$\rho_{\mathit{quant} \leftarrow \mathit{cache}}$:}
Sets the capacity contribution of set~$i$ to~0 when all lines in set~$i$ have
taint bit~0 after $\rho_{\mathit{cache} \leftarrow \mathit{spec}}$.  By
definition, untainted lines do not contribute to distinguishable configurations.

\emph{Monotonicity:} Both reduction steps are monotone with respect to the
product ordering: enlarging any component can only enlarge the set of feasible
paths, so the reduced result is also enlarged.

\emph{Termination:} The inner fixpoint of~$\rho$ terminates because $\Dcache$
has finite height $O(S \cdot A)$ and each application is non-increasing.
\end{proof}

\begin{theorem}[Composition Soundness]
\label{thm:composition-app}
Let $f, g$ be functions with contracts $(\tau_f, B_f)$ and $(\tau_g, B_g)$.
Under the independence condition,
$B_{f;g}(s) = B_f(s) + B_g(\tau_f(s))$
is a sound upper bound on the leakage of $f;g$.
\end{theorem}

\begin{proof}
By the sub-multiplicative counting argument:
$|\mathit{traces}_{f;g}| \le |\mathit{traces}_f| \times |\mathit{traces}_g|_{\tau_f(s)}$.
Taking logarithms yields additivity.  The independence condition ensures that
cache observations during~$g$ carry no information about $f$'s secrets beyond
what is captured by $\tau_f(s)$.  This instantiates
Smith~\cite{smith2009foundations}, Theorem~4.3, for the cache-state domain with
taint-restricted counting.
\end{proof}

\begin{lemma}[$\rho$ Monotonicity]
\label{lem:rho-mono-app}
For all $d_1 \sqsubseteq d_2 \in \Dspec \times \Dcache \times \Dquant$,
$\rho(d_1) \sqsubseteq \rho(d_2)$.
\end{lemma}

\begin{proof}
Follows from the monotonicity of projection and restriction on finite lattices.
\end{proof}

\begin{theorem}[Widening Correctness]
\label{thm:widening-app}
Let $\nabla$ denote the widening operator on $\Dspec \otimes \Dcache \otimes \Dquant$.
For any ascending chain $d_0 \sqsubseteq d_1 \sqsubseteq \cdots$ with
$d_{i+1} = F(d_i) \sqcup d_i$, the widened chain $d_0' = d_0$,
$d_{i+1}' = d_i' \nabla F(d_i')$ stabilizes in at most $h$ steps, where
$h = S \times A \times (W + 1)$ is the product lattice height.
\end{theorem}

\begin{proof}
The product lattice $\Dspec \otimes \Dcache \otimes \Dquant$ has finite
height bounded by $h = |\hat{S}| \times |\hat{C}| \times |\hat{Q}|$.
Each component widening operator guarantees strict ascent or stabilization.

\emph{$\nabla_{\text{spec}}$:}
The speculation domain has height $|\mathit{Loc}| \times (W + 2)$. Widening adds
at least one new reachable speculative state per step, ensuring termination
in at most $|\hat{S}|$ steps.

\emph{$\nabla_{\text{cache}}$:}
Cache-state widening widens age intervals monotonically. Since each age is
bounded by $[0, A]$, the maximum chain length per set is $A + 1$, giving
a total height of $S \times (A + 1)$.

\emph{$\nabla_{\text{quant}}$:}
Counting domain widening increases $h_i$ monotonically up to $\binom{A + L - 1}{A}$,
the maximum number of distinguishable configurations per set.

Since all three component widenings terminate and the product widening
applies them componentwise, the product chain stabilizes in at most
$h = S \times (A + 1) + |\mathit{Loc}| \times (W + 2)$ steps.
\end{proof}

\begin{lemma}[Taint Soundness]
\label{lem:taint-sound-app}
For any concrete execution trace $\pi$ of function $f$ with secret input $k$,
if instruction $i$ in $\pi$ accesses cache line $\ell$ and the access is
secret-dependent, then $\mathit{taint}(\alpha_{\text{cache}}(\pi), \ell) = 1$.
\end{lemma}

\begin{proof}
Secret-dependence of access $i$ means there exist secrets $k_1, k_2$ such that
$i$ accesses different cache lines under $k_1$ and $k_2$.  The abstract transfer
function for memory access sets $\mathit{taint}(\hat{c}, \ell) = 1$ whenever
the accessed address is not provably constant (i.e., the address abstract
value spans multiple cache lines).  Since the abstraction is sound, the
concrete access to $\ell$ under secret $k$ is included in the abstract
address range, and hence the taint bit is set.
\end{proof}

\section{Extended Evaluation Status}
\label{app:extended}

This appendix is intentionally brief. The repository already contains scripts
for richer evaluation tables, but the paper now defers those numerical details
until they are regenerated from clean, fully reproducible runs. The formal
definitions and proofs remain the stable appendix material for this version.

%% ===================================================================
%% REFERENCES
%% ===================================================================

\begin{thebibliography}{28}

\bibitem{cousot1977abstract}
P.~Cousot and R.~Cousot.
\newblock Abstract interpretation: A unified lattice model for static analysis
  of programs by construction or approximation of fixpoints.
\newblock In \emph{Proc.\ 4th ACM SIGACT--SIGPLAN Symposium on Principles of
  Programming Languages (POPL)}, pp.\ 238--252, 1977.

\bibitem{cousot1979systematic}
P.~Cousot and R.~Cousot.
\newblock Systematic design of program analysis frameworks.
\newblock In \emph{Proc.\ 6th ACM SIGACT--SIGPLAN Symposium on Principles of
  Programming Languages (POPL)}, pp.\ 269--282, 1979.

\bibitem{cousot1992abstract}
P.~Cousot and R.~Cousot.
\newblock Abstract interpretation frameworks.
\newblock \emph{Journal of Logic and Computation}, 2(4):511--547, 1992.

\bibitem{granger1992improving}
P.~Granger.
\newblock Improving the results of static analyses of programs by local
  decreasing iterations.
\newblock In \emph{Proc.\ Foundations of Software Technology and Theoretical
  Computer Science (FSTTCS)}, pp.\ 68--79, 1992.

\bibitem{doychev2015cacheaudit}
G.~Doychev, B.~K\"{o}pf, L.~Mauborgne, and J.~Reineke.
\newblock {CacheAudit}: A tool for the static analysis of cache side channels.
\newblock \emph{ACM Trans.\ Information and System Security (TISSEC)},
  18(1):4:1--4:32, 2015.

\bibitem{guarnieri2020spectector}
M.~Guarnieri, B.~K\"{o}pf, J.~F.~Morales, J.~Reineke, and
  A.~S\'{a}nchez.
\newblock {Spectector}: Principled detection of speculative information flows.
\newblock In \emph{Proc.\ IEEE Symposium on Security and Privacy (S\&P)},
  pp.\ 1--19, 2020.

\bibitem{daniel2020binsecrel}
L.~Daniel, S.~Bardin, and T.~Rezk.
\newblock {Binsec/Rel}: Efficient relational symbolic execution for
  constant-time at binary-level.
\newblock In \emph{Proc.\ IEEE Symposium on Security and Privacy (S\&P)},
  pp.\ 1003--1020, 2020.

\bibitem{gleissenthall2023leave}
Z.~Wang, G.~Mohr, K.~von Gleissenthall, J.~Reineke, and
  M.~Guarnieri.
\newblock Specification and verification of side-channel security for
  open-source processors via leakage contracts.
\newblock In \emph{Proc.\ ACM Conference on Computer and Communications
  Security (CCS)}, pp.\ 2411--2425, 2023.
\newblock (Distinguished Paper Award.)

\bibitem{smith2009foundations}
G.~Smith.
\newblock On the foundations of quantitative information flow.
\newblock In \emph{Proc.\ International Conference on Foundations of Software
  Science and Computational Structures (FoSSaCS)}, pp.\ 288--302, 2009.

\bibitem{kopf2010approximation}
B.~K\"{o}pf and A.~Rybalchenko.
\newblock Approximation and randomization for quantitative information-flow
  analysis.
\newblock In \emph{Proc.\ IEEE Computer Security Foundations Symposium (CSF)},
  pp.\ 3--14, 2010.

\bibitem{dwork2006differential}
C.~Dwork.
\newblock Differential privacy.
\newblock In \emph{Proc.\ International Colloquium on Automata, Languages, and
  Programming (ICALP)}, pp.\ 1--12, 2006.

\bibitem{myers1999jflow}
A.~C.~Myers.
\newblock {JFlow}: Practical mostly-static information flow control.
\newblock In \emph{Proc.\ 26th ACM SIGPLAN--SIGACT Symposium on Principles of
  Programming Languages (POPL)}, pp.\ 228--241, 1999.

\bibitem{enck2010taintdroid}
W.~Enck, P.~Gilbert, B.-G.~Chun, L.~P.~Cox, J.~Jung, P.~McDaniel, and
  A.~N.~Sheth.
\newblock {TaintDroid}: An information-flow tracking system for realtime privacy
  monitoring on smartphones.
\newblock In \emph{Proc.\ 9th USENIX Symposium on Operating Systems Design and
  Implementation (OSDI)}, pp.\ 393--407, 2010.

\bibitem{reineke2007timing}
J.~Reineke, D.~Grund, C.~Berg, and R.~Wilhelm.
\newblock Timing predictability of cache replacement policies.
\newblock \emph{Real-Time Systems}, 37(2):99--122, 2007.

\bibitem{ge2018survey}
Q.~Ge, Y.~Yarom, D.~Cock, and G.~Heiser.
\newblock A survey of microarchitectural timing attacks and countermeasures on
  contemporary hardware.
\newblock \emph{Journal of Cryptographic Engineering}, 8(1):1--27, 2018.

\bibitem{gruss2016flush}
D.~Gruss, C.~Maurice, K.~Wagner, and S.~Mangard.
\newblock {Flush+Flush}: A fast and stealthy cache attack.
\newblock In \emph{Proc.\ Conference on Detection of Intrusions and Malware,
  and Vulnerability Assessment (DIMVA)}, pp.\ 279--299, 2016.

\bibitem{yarom2014flush}
Y.~Yarom and K.~Falkner.
\newblock {FLUSH+RELOAD}: A high resolution, low noise, {L3} cache
  side-channel attack.
\newblock In \emph{Proc.\ 23rd USENIX Security Symposium}, pp.\ 719--732,
  2014.

\bibitem{kocher2019spectre}
P.~Kocher, J.~Horn, A.~Fogh, D.~Genkin, D.~Gruss, W.~Haas, M.~Hamburg,
  M.~Lipp, S.~Mangard, T.~Prescher, M.~Schwarz, and Y.~Yarom.
\newblock Spectre attacks: Exploiting speculative execution.
\newblock In \emph{Proc.\ IEEE Symposium on Security and Privacy (S\&P)},
  pp.\ 1--19, 2019.

\bibitem{alvim2020science}
M.~S.~Alvim, K.~Chatzikokolakis, A.~McIver, C.~Morgan, C.~Palamidessi, and
  G.~Smith.
\newblock \emph{The Science of Quantitative Information Flow}.
\newblock Springer, 2020.

\bibitem{almeida2016verifying}
J.~B.~Almeida, M.~Barbosa, G.~Barthe, F.~Dupressoir, and M.~Emmi.
\newblock Verifying constant-time implementations.
\newblock In \emph{Proc.\ 25th USENIX Security Symposium},
  pp.\ 53--70, 2016.

\bibitem{cauligi2020constant}
S.~Cauligi, C.~Disselkoen, K.~von Gleissenthall, D.~Tullsen,
  D.~Stefan, T.~Rezk, and G.~Barthe.
\newblock Constant-time foundations for the new {Spectre} era.
\newblock In \emph{Proc.\ ACM SIGPLAN Conference on Programming Language
  Design and Implementation (PLDI)}, pp.\ 913--926, 2020.

\bibitem{wang2017cached}
S.~Wang, P.~Wang, X.~Liu, D.~Zhang, and A.~Wu.
\newblock {CacheD}: Identifying cache-based timing channels in production
  software.
\newblock In \emph{Proc.\ 26th USENIX Security Symposium},
  pp.\ 235--252, 2017.

\bibitem{simon2018you}
L.~Simon, D.~Chisnall, and R.~Anderson.
\newblock What you get is what you {C}: Controlling side effects in mainstream
  {C} compilers.
\newblock In \emph{Proc.\ IEEE European Symposium on Security and Privacy
  (EuroS\&P)}, pp.\ 1--15, 2018.

\bibitem{bernstein2005poly1305}
D.~J.~Bernstein.
\newblock The {Poly1305-AES} message-authentication code.
\newblock In \emph{Proc.\ Fast Software Encryption (FSE)},
  pp.\ 32--49, 2005.

\bibitem{brumley2011remote}
B.~B.~Brumley and N.~Tuveri.
\newblock Remote timing attacks are still practical.
\newblock In \emph{Proc.\ European Symposium on Research in Computer Security
  (ESORICS)}, pp.\ 355--371, 2011.

\bibitem{osvik2006cache}
D.~A.~Osvik, A.~Shamir, and E.~Tromer.
\newblock Cache attacks and countermeasures: The case of {AES}.
\newblock In \emph{Proc.\ Cryptographers' Track at the RSA Conference
  (CT-RSA)}, pp.\ 1--20, 2006.

\bibitem{doychev2017rigorous}
G.~Doychev and B.~K\"{o}pf.
\newblock Rigorous analysis of software countermeasures against cache attacks.
\newblock In \emph{Proc.\ ACM SIGPLAN Conference on Programming Language
  Design and Implementation (PLDI)}, pp.\ 406--421, 2017.

\bibitem{barthe2020formal}
G.~Barthe, S.~Blazy, B.~Gr\'{e}goire, R.~Hutin, V.~Laporte,
  D.~Pichardie, and A.~Trieu.
\newblock Formal verification of a constant-time preserving {C} compiler.
\newblock \emph{Proc.\ ACM on Programming Languages}, 4(POPL):7:1--7:30, 2020.

\bibitem{bourdoncle1993efficient}
F.~Bourdoncle.
\newblock Efficient chaotic iteration strategies with widenings.
\newblock In \emph{Proc.\ International Conference on Formal Methods in
  Programming and Their Applications}, pp.\ 128--141, 1993.

\bibitem{watt2019ct}
C.~Watt, J.~Renner, N.~Popescu, S.~Cauligi, and D.~Stefan.
\newblock {CT-Wasm}: Type-driven secure cryptography for the web ecosystem.
\newblock \emph{Proc.\ ACM on Programming Languages}, 3(POPL):77:1--77:29,
  2019.

\bibitem{wichelmann2018microwalk}
J.~Wichelmann, A.~Moghimi, T.~Eisenbarth, and B.~Sunar.
\newblock {MicroWalk}: A framework for finding side channels in binaries.
\newblock In \emph{Proc.\ 34th Annual Computer Security Applications
  Conference (ACSAC)}, pp.\ 161--173, 2018.

\bibitem{bond2017vale}
B.~Bond, C.~Hawblitzel, M.~Kapritsos, K.~R.~M.~Leino, J.~R.~Lorch,
  B.~Parno, A.~Rane, S.~Setty, and L.~Thompson.
\newblock {Vale}: Verifying high-performance cryptographic assembly code.
\newblock In \emph{Proc.\ 26th USENIX Security Symposium},
  pp.\ 917--934, 2017.

\bibitem{cauligi2019fact}
S.~Cauligi, G.~Soeller, B.~Johannesmeyer, F.~Brown, R.~S.~Wahby,
  J.~Renner, B.~Gr\'{e}goire, G.~Barthe, R.~Jhala, and D.~Stefan.
\newblock {FaCT}: A {DSL} for timing-sensitive computation.
\newblock In \emph{Proc.\ ACM SIGPLAN Conference on Programming Language
  Design and Implementation (PLDI)}, pp.\ 174--189, 2019.

\bibitem{bearssl}
T.~Pornin.
\newblock {BearSSL}: A smaller {SSL/TLS} library.
\newblock \url{https://bearssl.org/}, 2016--2023.

\bibitem{signal2016}
Open Whisper Systems.
\newblock Signal Protocol library ({C} implementation).
\newblock \url{https://github.com/signalapp/libsignal-protocol-c}, 2016.

\end{thebibliography}

\end{document}
